\chapter{Methodology}

\section{Principles}

The protocol distinguishes:
\begin{description}
  \item[R1:] explicitly published information;
  \item[R2:] reconstruction to fill a gap;
  \item[A1:] adaptation required for crypto;
  \item[E1:] extension with a Transformer.
\end{description}

\section{Pipeline}

\begin{enumerate}
  \item build the daily URL and download the ZIP;
  \item compute the SHA-256 and record it in a manifest;
  \item temporarily extract the CSV;
  \item process chunks of one million rows;
  \item reduce the updates to the 100 ms grid;
  \item produce Parquet partitioned by symbol and date;
  \item delete the temporary CSV;
  \item create windows without crossing splits;
  \item fit the MLE and calibrate ranges on the training set only;
  \item generate synthetic paths and train the estimators;
  \item select hyperparameters on the validation set;
  \item open the test set only once.
\end{enumerate}

\section{Simulation by thinning}
\label{sec:sim-thinning}

Given the current state, an upper bound on the sum of the intensities is computed. An
exponential waiting time is drawn; the excitation decays; the candidate is accepted
with probability proportional to the new intensity \citep{ogata1981}. The type is
drawn with probabilities proportional to \(\lambda_1,\lambda_2\). After the event, the
corresponding column of the \(\alpha\) matrix is added. Each path starts with an
empty history, the same assumption as in the likelihood.

\section{Constrained MLE}

Positive parameters are obtained through softplus. Instead of optimizing \(\beta\)
freely, a positive margin \(d\) is used:
\[
(\mu,\alpha_s,\alpha_c,d)=\softplus(z),\qquad
\beta=\alpha_s+\alpha_c+d.
\]
L-BFGS-B is started from five points: one based on the empirical event rate and four
random perturbations in the unconstrained space. The MLE never receives the true
parameter as initialization. The gradient is analytic, and the exponential
recursions are solved as bidiagonal systems. Convergence, iterations, log-likelihood
and time are recorded.

\section{Neural protocol}

The target is \((\mu,\alpha_s,\alpha_c,\beta)\), and the loss is the MSE on the
original scale. The input for each event is \((\log(1+\Delta t),\,s)\), with \(s=+1\)
for an up-move and \(s=-1\) for a down-move; Lee uses inter-arrival times and types
1/2, so this transformation is a reconstruction (R2). The two networks share the
output head: softplus produces \(\mu,\alpha_s,\alpha_c,d>0\) and
\(\beta=\alpha_s+\alpha_c+d\), which guarantees \(g<1\) by construction, as in Lee's
robust version.

The frozen synthetic experiment contains 75,000 training sequences, 10,000 validation
sequences and 15,000 test sequences, all of exactly 1,000 events. The parameters are
drawn with \(\mu\in[0.2,\,0.6]\), \(\beta\in[0.8,\,2.0]\), self-excited fraction
\(\alpha_s/(\alpha_s+\alpha_c)\in[0.2,\,0.8]\) and \(g\) balanced across three
strata: low \([0.10,\,0.35)\), intermediate \([0.35,\,0.70)\) and near-critical
\([0.70,\,0.95]\). Table~\ref{tab:treino} summarizes the training.

\begin{table}[htbp]
\centering
\caption{Training configuration and cost (five seeds per network, Apple Silicon Mac
with MPS).}
\label{tab:treino}
\small
\begin{tabular}{lrr}
\toprule
 & LSTM (R1/R2) & Transformer (E1)\\
\midrule
Architecture & 2 LSTM $\times$ 12 + dense & encoder, 2 layers, $d=32$\\
Heads / feed-forward & -- & 4 / 128\\
Trainable parameters & 2,068 & 25,700\\
Optimizer & Adam, $10^{-3}$ & Adam, $10^{-3}$\\
Effective batch & 64 & 8 (physical 4 with accumulation)\\
Maximum epochs / patience & 300 / 30 & 300 / 30\\
Best epoch (per seed) & 165, 164, 193, 239, 299 & 73, 42, 38, 60, 44\\
Training time per seed & 5.4--14.8 h & 16.6--28.8 h\\
Total time (5 seeds) & 41 h & 110 h\\
\bottomrule
\end{tabular}
\end{table}

The Transformer's tensor flow starts from \([B,L,2]\), with \(L=1000\). A linear
projection produces \([B,L,32]\), to which the positional encoding is added. In each
layer, the four heads operate in subspaces of dimension eight; the feed-forward block
expands 32 to 128 and returns to 32. Temporal averaging converts \([B,L,32]\) into
\([B,32]\). The memory cost of attention grows quadratically with \(L\), which is why
its batch is smaller than the LSTM's. This difference in batch size, together with the
absence of a hyperparameter search for either network, is a limitation of the
comparison.

\section{The paper's fixed point (E2)}

To compare directly with the sampling-distribution table of \citet{lee2023}, 2,000
paths of 1,000 events were simulated at the published point
\((\mu,\alpha_s,\alpha_c,\beta)=(0.3,\,0.4,\,0.7,\,1.5)\), which lies inside the
training prior. The networks were applied without retraining, and the MLE was fitted
in two versions: the protocol's heuristic and the initialization at the true value
used in the paper.

\section{Empirical protocol}
\label{sec:empirical-protocol}

On real data there is no true parameter, so the empirical stage asks three narrower
questions: whether the Hawkes model describes the windows well enough to be estimated, whether
a network can reproduce the MLE on real windows, and whether the estimated parameters give a
volatility consistent with the realized one. The protocol was fixed in a decision record
(ADR~008) before the test block was opened. The windows have 1,000 events and are split
chronologically by day: calibration (15--22 January 2024, 1,063 windows), validation (23--25
January, 476 windows) and test (26--28 January, 274 windows); the October 2023 week is a
robustness window that is never used for training or selection.

\paragraph{Goodness of fit.} With the MLE of each calibration window, the time-rescaled
residuals of each dimension are tested against \(\mathrm{Exp}(1)\) with Kolmogorov--Smirnov and
for serial dependence with Ljung--Box at 10 lags \citep{ogata1988,brown2002}, and the Hawkes
model is compared with a homogeneous bivariate Poisson process by likelihood ratio and AIC.
Since the times lie on a 100~ms grid, two controls separate the effect of the grid from
misspecification: times jittered uniformly inside their cell with the MLE refitted, and paths
simulated from each fitted model, either in continuous time or mapped to the grid in the same
way as the data. The rejection rate of the simulated controls is what a correct model would
produce.

\paragraph{Networks on real windows.} Each window is labeled with its MLE (five heuristic
starts). Both architectures are trained on the calibration labels with the loss, batch size,
maximum of 300 epochs and patience of 30 of the synthetic protocol, in two variants: from
scratch (Adam, \(10^{-3}\)) and fine-tuned from the synthetic weights of the same seed (Adam,
\(10^{-4}\)), each with the five protocol seeds. Two references need no training: the synthetic
networks as they are, and the mean of the training labels, which is the best constant
predictor. The selected variant is the one whose five-seed ensemble has the lowest macro MSE
on validation. The test block is labeled and read once, after the selection; the stored
weights are evaluated without retraining, and every variant is reported. The machine budget
for this stage was capped at 30 hours.

\paragraph{Volatility.} The long-run volatility implied by the estimated parameters
(Section~\ref{sec:real-volatility}) is compared window by window with the realized volatility
of the log mid-price sampled every \(\Delta\in\{0.1, 1, 5, 10, 30\}\) seconds, using the median
ratio and the Spearman correlation across windows, with \(\Delta=10\)~s as the main horizon,
for the MLE and for the selected network.

\section{Metrics and tests}

The reported metrics are per-parameter and aggregate MSE, median absolute error,
bias, validity, branching-ratio error, cost and MLE convergence. Because the
parameters have different scales, the \(R^2\) of each parameter relative to the prior
mean (the trivial estimator) is also reported. The comparison between networks uses
two sources of variation: across seeds (Welch and Mann--Whitney tests with five
seeds), which measures training variability \citep{bouthillier2021}, and across test
paths (paired difference with a \emph{bootstrap} CI and Wilcoxon test), which measures
the accuracy of the already trained models \citep{efron1993,diebold1995}.
